% part: intuitionistic-logic
% chapter: introduction
% section: constructive-reasoning

\documentclass[../../../include/open-logic-chapter]{subfiles}

\begin{document}

\olfileid[tr]{int}{int}{cr}

\section{Yapıcı Akıl Yürütme}

Klasik mantığın kiplik işleçleriyle ya da ikinci dereceden niceleyicilerle
genişletilmesinden farklı olarak, sezgici mantık klasik mantığı kısıtladığı için
``klasik olmayan'' bir mantıktır. Klasik mantık çeşitli bakımlardan
\emph{yapıcı değildir}. Sezgici mantık, yapıcı matematiğin bir türüne özgü,
daha ``yapıcı'' bir akıl yürütme türünü yakalamayı amaçlar. Aşağıdaki örnekler,
bunun altında yatan güdülerden bazılarını açıklamaya yarayabilir.

Birinin, şu özelliğe sahip bir $n$ doğal sayısı belirlediğini ileri sürdüğünü
varsayın: $n$ çiftse Riemann hipotezi doğru, $n$ tekse Riemann hipotezi
yanlıştır. Müthiş haber!{} Riemann hipotezinin doğru olup olmadığı matematiğin
büyük açık sorularından biridir; oysa bu kişi sorunu bir hesaplama problemine,
yani belirli bir sayının çift olup olmadığını belirlemeye indirgemiş
görünmektedir.

$n$'nin sihirli değeri nedir? Şöyle betimliyor: $n$, Riemann hipotezi doğruysa
$2$'ye, değilse $3$'e eşit olan doğal sayıdır.

Öfkelenip paranızı geri istersiniz. Klasik bakımdan yukarıdaki betimleme
gerçekten de $n$'nin tek bir değerini belirler; ama sizin asıl istediğiniz,
\emph{açıkça} verilen bir $n$ değeridir.

Başka, belki daha az yapay bir örnek olarak şu soruyu ele alın. İrrasyonel bir
sayının rasyonel bir kuvvetini alıp rasyonel bir sonuç elde etmenin mümkün
olduğunu biliyoruz. Örneğin, $\sqrt{2}^2 = 2$. Daha az açık olan, irrasyonel bir
sayının \emph{irrasyonel} bir kuvvetini alıp rasyonel bir sonuç elde etmenin
mümkün olup olmadığıdır. Aşağıdaki teorem buna olumlu yanıt verir:

\begin{thm}
Öyle irrasyonel $a$ ve $b$ sayıları vardır ki $a^b$ rasyoneldir.
\end{thm}

\begin{proof}
$\sqrt{2}^{\sqrt{2}}$ sayısını ele alın. Bu sayı rasyonelse işimiz bitmiştir:
$a = b = \sqrt{2}$ alabiliriz. Değilse irrasyoneldir. O zaman
\[
(\sqrt{2}^{\sqrt{2}})^{\sqrt{2}} = \sqrt{2}^{\sqrt{2} \cdot
  \sqrt{2}} = \sqrt{2}^2 = 2,
\]
olur ve bu sayı rasyoneldir. Dolayısıyla bu durumda $a$'yı
$\sqrt{2}^{\sqrt{2}}$, $b$'yi de~$\sqrt 2$ alalım.
\end{proof}

Bu, geçerli bir ispat sayılır mı? Çoğu matematikçi öyle olduğunu düşünür.
Ancak burada da biraz doyurucu olmayan bir şey vardır: belirli bir özelliğe
sahip bir gerçek sayı çiftinin varlığını, bunun \emph{hangi} sayı çifti
olduğunu söyleyemeden ispatladık. Aynı sonuç, $a$, $b$ çifti ispatta
\emph{verilecek} biçimde de ispatlanabilir: $a = \sqrt{3}$ ve $b = \log_3 4$
alın. O zaman
\[
a^b = \sqrt{3}^{\log_3 4} = 3^{1/2 \cdot \log_3 4} = (3^{\log_3
  4})^{1/2} = 4^{1/2}= 2,
\]
elde edilir; çünkü $3^{\log_3 x} = x$.

Sezgici mantık, ilk ispattaki gibi hamlelere izin verilmeyen bir akıl yürütme
türünü yakalamak üzere tasarlanmıştır. Bir $x$'in~$!A(x)$ koşulunu sağladığını
göstererek varlığını ispatlamak, ikinci ispattaki gibi belirli bir~$x$ ile onun $!A$
koşulunu sağladığının bir ispatını vermeniz gerektiği anlamına gelir. $!A$ ya
da $!B$'nin geçerli olduğunu ispatlamak da bunlardan birini ispatlayabilmenizi
gerektirir.

Biçimsel olarak sezgici mantık, klasik mantığa ait !!a{derivation} sistemini
belirli bir biçimde kısıtlayınca elde edilen şeydir. Matematiksel bakımdan
bunlar yalnızca biçimsel tümdengelim sistemleridir; ama daha önce belirtildiği
gibi belirli bir matematiksel akıl yürütme türünü yakalamaları amaçlanır. Bu,
matematik hakkındaki belirli bir felsefi görüşün (Brouwer'ın sezgiciliği gibi)
gerekçelendirdiği akıl yürütme türü olarak görülebilir; daha ``somut'' ve
doyurucu bir matematiksel akıl yürütme türü (Bishop'ın yapıcılığı doğrultusunda)
sayılabilir; biçimsel betimlemenin biçimsel olmayan güdüyü yakalayıp
yakaladığı da tartışılabilir. Fakat hangi felsefi görüşü benimsersek
benimseyelim, sezgici mantığı biçimsel olarak sunulmuş bir mantık diye
inceleyebiliriz; her ne sebeple olursa olsun, birçok matematiksel mantıkçı bunu
yapmayı ilginç bulur.

\end{document}
